\chapter{Objetivos}\label{cap:objetivos}

\section{Objetivo general}

El objetivo general del TFM, según los objetivos oficiales propuestos para el TFE en la universidad, es \textbf{diseñar y desplegar una configuración de metadatos en OpenMetadata alineada con el estándar DCAT-AP}, que permita \textbf{mejorar la interoperabilidad y la gestión semántica de catálogos de datos} en entornos Big Data y de computación en la nube. En este repositorio, ese perfil DCAT-AP se concreta en su aplicación nacional DCAT-AP-ES.

En términos implementados, este objetivo se materializa en una Plataforma de Gobierno del Dato capaz de descubrir activos técnicos, enriquecerlos con metadatos funcionales, exportar un catálogo JSON-LD y validar su conformidad mediante SHACL. La validación se apoya en dos servicios PostgreSQL ingeridos en OpenMetadata, lo que permite demostrar que el gobierno no depende de una única fuente ni de nombres de tabla globalmente únicos.

\section{Objetivos específicos}

Los objetivos específicos se formulan de forma verificable:

\begin{enumerate}
  \item Analizar DCAT-AP y su concreción DCAT-AP-ES, incluyendo clases, propiedades, vocabularios controlados y validación SHACL.
  \item Mapear las entidades de OpenMetadata con las clases principales de DCAT-AP-ES.
  \item Configurar en OpenMetadata los metadatos mínimos necesarios para representar datasets publicables.
  \item Implementar un \emph{workflow} reproducible de sincronización de metadatos desde activos técnicos reales.
  \item Ejecutar doble ingesta PostgreSQL para validar gobierno multi-fuente con tablas homónimas.
  \item Exportar el catálogo gobernado en JSON-LD compatible con DCAT-AP-ES.
  \item Validar la salida con SHACL y evaluar beneficios, limitaciones, idempotencia y mantenibilidad.
  \item Proporcionar una consola web operativa que ejecute el mismo flujo sin duplicar reglas de negocio.
\end{enumerate}

\section{Cambio de alcance respecto a CKAN}

La descripción oficial del TFE planteaba implementar un \emph{pipeline} de ingesta o sincronización de metadatos (\emph{harvesting}) desde una fuente externa, y proponía CKAN como ejemplo. Durante el desarrollo se optó por no emplear CKAN como origen operativo principal. La razón es de fondo: CKAN es ante todo un catálogo de publicación e intercambio de metadatos ya elaborados, de modo que cosechar metadatos de un catálogo externo no demuestra la parte central del trabajo, que es generar, gobernar y validar metadatos a partir de activos técnicos reproducibles. La formulación literal de la descripción no resultaba, por tanto, del todo apropiada para el objetivo de fondo del TFM. El cambio de origen queda suficientemente justificado en términos de funcionalidad y de objetivos, porque permite demostrar el ciclo completo de gobierno que el trabajo persigue sin alterar lo esencial de lo planteado: alinear OpenMetadata con DCAT-AP y validar la interoperabilidad del catálogo resultante.

Conviene además separar dos conceptos que no son equivalentes. La \emph{gestión del dato} (\emph{data management}) abarca las funciones operativas que producen y mantienen los datos (arquitectura, modelado, almacenamiento, integración o calidad). El \emph{gobierno del dato} (\emph{data governance}) es, en cambio, la función de supervisión que define políticas, roles, responsabilidades y controles sobre esas funciones; en el marco DAMA-DMBOK es el núcleo que coordina y da autoridad al resto de la gestión, no una tarea operativa más~\cite{damaDmbok,datosGobDama2021}. Este TFM se sitúa en la función de gobierno (decidir qué se publica, con qué metadatos y bajo qué reglas), apoyándose en funciones de gestión como la catalogación y la integración sin pretender cubrirlas todas. CKAN se mantiene como posible destino o punto de federación futura, no como fuente canónica del caso de uso de validación.
